% !TEX program = xelatex
% ============================================================
% 2026 全国大学生数学建模竞赛（CUMCM）电子版论文模板
% 编译：XeLaTeX / latexmk
% ------------------------------------------------------------
% 1. 比赛写作主要修改本 main.tex；
% 2. 图片全部放 figs/；源码和支撑材料放 support/；
% 3. 【...】为模板占位内容，正式提交前必须全部替换或删除；
% 4. 正式提交前运行：python tools/preflight.py
% ============================================================

\documentclass{cumcm2026}
% ============================================================
% 用户导言区：自己的宏、额外宏包统一放这里。
% 不建议修改 cumcm2026.cls，避免比赛中误伤版式。
% ============================================================
\newcommand{\R}{\mathbb{R}}
\newcommand{\E}{\mathbb{E}}
\newcommand{\Var}{\operatorname{Var}}
\newcommand{\Cov}{\operatorname{Cov}}

% ============================================================
% BibTeX 参考文献管理（移植自 latex 模板）
% ------------------------------------------------------------
% 使用方法：
%   1. 全部真实文献统一维护在项目根目录的 ref.bib 中；
%   2. 正文中用 \cite{key} / \citep{key} / \citet{key} 引用；
%   3. 引用样式由根目录 gbt7714-numerical.bst 控制（GB/T 7714 数字式）。
% 编译：latexmk 会自动按 xelatex→bibtex→xelatex×2 的顺序处理；
% 手动则执行：xelatex main && bibtex main && xelatex main && xelatex main。
% ============================================================
\usepackage[numbers,sort&compress]{natbib}
% 让“参考文献”继续作为带编号的正文小节（与 cumcm2026.cls 的章节风格一致）
\renewcommand{\bibsection}{\section{\refname}}


\title{基于洪水假设的多模态分析和解析}

\begin{document}

% ============================================================
% 摘要专用页：电子版第一页必须是摘要页，且摘要原则上不超过一页
% ============================================================
\maketitle
\begin{abstract}

本文针对题目核心对象的预测/评价/优化/决策问题展开研究。首先对题目提供的数据进行缺失值处理、异常值识别与标准化，并结合题意提取关键变量1、关键变量2和关键变量3等主要影响因素。在此基础上，分别针对各子问题建立相应的数学模型。

针对问题一，本文采用模型/方法名\textbf{称刻画变量}

之间的关系，得到最重要的数值结果或规律；针对问题二，在问题一基础上进一步考虑新增约束或影响因素，建立模型名称，求得最优方案为核心方案，对应目标值为核心数值；针对问题三，引入算法或情景分析方法对复杂情形进行求解，并通过敏感性分析/误差分析/蒙特卡洛模拟检验模型的稳定性。结果表明，当关键参数在范围内变化时，主要结论保持稳定，说明模型具有较好的鲁棒性。

综合来看，本文模型能够较好地描述实际机制，具有结构清晰、可解释性强、求解效率高等优点，可为实际应用场景提供定量决策依据。

\keywords{数学建模\quad 关键词2\quad 关键词3\quad 关键词4}
\end{abstract}

% 此标签供提交前检查脚本验证“摘要只占第一页”。不要删除。
\phantomsection\label{cumcm:body-start}

% CUMCM 正文不要目录，直接开始正文。

% ============================================================
% 一、问题重述
% ============================================================
\section{问题重述}

\subsection{问题背景}

微网是促进分布式光伏本地消纳、降低外部电网供电压力并提高能源利用效率的重要形式。非孤岛式微网与外部电网相连，既可在光伏出力不足或小区负载较高时从外网购电，也可在电价较低或光伏有富余时为储能设备充电，并在电价较高时释放电能。因此，微网运行需要在满足小区负载和储能安全要求的前提下，合理协调光伏、储能与外部电网之间的电能流动，尽可能降低购电成本。

在实际运行中，小区负载和光伏出力具有明显的时段变化特征，外部电网电价也可能保持固定曲线或发生实时波动。若仅依据当前状态进行购电，容易错失低价充电和高价放电机会；若储能调度不当，还可能造成过充、过放、跨日状态中断或紧急购电费用增加。因此，需要根据负载、光伏、电价和储能状态等数据，在连续时间尺度上制定购电量和充放电量，使微网在满足供电要求的同时获得经济、可行且稳定的调度方案。

% ---- 参考文献引用示例（BibTeX 用法）：写作时把用到的文献 key 填进 \cite{...}，条目统一写进根目录 ref.bib 即可 ----
本文在建模中将参考数学建模专著\cite{司守奎2011数学建模算法与应用}、MATLAB 建模工具书\cite{卓金武2011MATLAB}以及多准则评价方法综述\cite{behzadian2012topsis}等文献，后续各模型所用方法均按此方式标注引用。

\subsection{基本问题}
在实际运行中，小区负载和光伏功率具有明显的时段变化特征，外部电网电价也可能发生实时波动。为逐步优化多动态因素下微网购电量方案，题目循序渐进地给出四个小问题和相关的约束条件。附录1给出了储能设备的相关约束条件：储能设备的最大容量为12000 kWh，最大充放电功率为5000 kW，2025年1月1日0:00的初始储电量为6000 kWh。为避免过度充电或放电，储能运行过程中的储电量必须保持在1200至10800 kWh之间，充放电效率均为90\%。各时段还需满足功率平衡、充放电功率上限和储能状态递推约束。微网提供的电能不得低于小区负载；在储能状态需要跨日延续时，下一日0:00的储电量应与前一日24:00的储电量保持一致。

附件1给出了某天电价、小区负载、光伏发电预测功率全天144个时间段的数值；附件2给出了
2025年1月1日至12月31日每天小区负载和光伏实际功率按按10分钟间隔记录的数据；附件3给出了2025年每天0:00、6:00、12:00和18:00预测的未来24小时光伏功率。基于这些数据集，完成以下问题：

1.使用附件1给定的电价、小区负载、光伏发电预测功率数据，分析小区负载、光伏功率和电价之间的关系，并建立优化模型得到全天最优购电量的方案，并给出指定时段购电量和储能设备的充放电量。

2.在问题一固定的电价的基础上，考虑附件2中小区负载、光伏发电实际功率随日期的变化，建立时序预测模型，在电能低于负载时额外5倍的紧急购电费的惩罚下，构建优化模型，给出全年每日最优购电量方案。

3.基于问题二的固定的电价和小区负载，考虑附件3中光伏发电功率和附件3中各时间段未来24小时光伏功率预测数据，结合紧急购电规则和计划-调整规则建立数学模型，得出全每天0:00的最优计划购电方案和在其他预报时刻的最优调整购电方案。

4.在问题二和三的基础上，考虑电价的实时波动，改进优化模型，重新计算问题二和三，分别给出全年最优购电量方案。
\begin{enumerate}[label=\textbf{问题\arabic*：}]
  \item 根据已知数据，分析【对象A】与【指标B】之间的关系，并建立模型得到【目标结果】；
  \item 在问题一的基础上考虑【新增因素/约束】，求得【预测结果、评价结果或最优方案】；
  \item 针对【更复杂情形】，建立综合模型，并对结果的稳定性和实际意义进行分析。
\end{enumerate}

% ============================================================
% 二、问题分析
% ============================================================
\section{问题分析}

问题分析的重点不是重复题意，而是解释“为什么采用后面的模型”。建议从数据特征、变量关系、约束条件和求解目标四个方面展开。

\subsection{问题一分析}
【示例】问题一的核心是刻画【自变量】对【因变量】的影响。由于数据具有【连续性/离散性/时间性/非线性】等特征，因此首先进行【相关性分析、特征工程或数据变换】，随后采用【模型名称】完成【预测/评价/拟合】。该模型的输出与题目要求的【指标】直接对应。

\subsection{问题二分析}
【示例】问题二在问题一基础上新增【成本、容量、时间、风险等】约束，因此不能只追求单一指标最优。可将其转化为带约束的【单目标/多目标】优化问题，以【目标函数】为优化目标，同时满足【主要约束】。

\subsection{问题三分析}
【示例】问题三需要进一步考察模型在不同情景下的适用性。为此，在前述模型基础上改变【关键参数】，利用【灵敏度分析/情景模拟/蒙特卡洛方法】比较输出变化，并据此评价模型稳定性与推广能力。

% ============================================================
% 三、模型假设
% ============================================================
\section{模型假设}
根据题意并结合实际情况，作如下合理假设。正式论文中只保留真正服务于模型的假设。
\begin{enumerate}[label=\textbf{假设\arabic*：}]
  \item 题目所给数据真实可靠，除已识别异常值外不存在系统性记录错误；
  \item 在研究时间范围内，题目未明确考虑的外部环境因素保持相对稳定；
  \item 同一指标采用统一统计口径，不同样本之间具有可比性；
  \item 优化模型中的所有决策变量均满足题目给出的实际可行域约束。
\end{enumerate}

% ============================================================
% 四、符号说明
% ============================================================
\section{符号说明}
为便于后文表述，定义主要符号如\cref{tab:symbols}所示。只建议列高频核心符号，临时变量可在首次出现时说明。

\begin{table}[htbp]
  \centering
  \caption{主要符号及含义}
  \label{tab:symbols}
  \begin{tabularx}{0.90\textwidth}{cXc}
    \toprule
    符号 & 含义 & 单位 \\
    \midrule
    $x_i$ & 第 $i$ 个样本的【特征/决策变量】 & 【单位】 \\
    $y_i$ & 第 $i$ 个样本的【观测值/目标变量】 & 【单位】 \\
    $w_j$ & 第 $j$ 个指标的权重 & -- \\
    $n$ & 样本总数 & -- \\
    $F(\bm{x})$ & 优化模型目标函数 & -- \\
    \bottomrule
  \end{tabularx}
\end{table}

% ============================================================
% 五、数据与预处理
% ============================================================
\section{数据与预处理}

\subsection{数据来源与说明}
【示例】题目直接给出的原始数据用于【用途】；另通过【公开数据库/政府网站/文献】获得【外部数据】用于【用途】。自行查阅并实际使用的数据应同时放入支撑材料；若确需引用外部资料，在正文对应位置用 \verb|\cite{key}| 标注，其中 key 为根目录 ref.bib 中该条目的键名（用法见“问题重述”一节中的示例）。若没有使用题面之外的数据，直接说明即可。

\subsection{缺失值与异常值处理}
【示例】首先统计各变量缺失比例。对于少量随机缺失值，采用【均值/中位数/插值】补全；对于明显异常且无法由实际机制解释的数据点，依据【3$\sigma$/箱线图/IQR/业务规则】进行识别与处理。

\subsection{数据标准化}
若不同指标量纲差异较大，可采用极差标准化：
\begin{equation}
  z_{ij}=\frac{x_{ij}-\min_i x_{ij}}{\max_i x_{ij}-\min_i x_{ij}}.
  \label{eq:normalization}
\end{equation}
经处理后，各指标处于统一尺度，便于后续模型计算与比较。

% ============================================================
% 六、模型建立与求解
% ============================================================
\section{模型建立与求解}

\subsection{问题一模型的建立}
【示例】设目标变量为 $y$，影响因素为 $x_1,x_2,\ldots,x_p$。若采用多元线性回归，可建立
\begin{equation}
  y=\beta_0+\sum_{j=1}^{p}\beta_j x_j+\varepsilon,
  \label{eq:q1-model}
\end{equation}
其中 $\beta_j$ 为待估参数，$\varepsilon$ 为随机误差项。若实际数据呈明显非线性，则应将本示例替换为更合适的模型。

\subsection{问题一模型的求解与结果}
利用【Python/MATLAB/R】对模型进行求解，得到关键指标如\cref{tab:q1-result}所示。

\begin{table}[htbp]
  \centering
  \caption{问题一结果示例}
  \label{tab:q1-result}
  \begin{tabular}{cccc}
    \toprule
    指标 & 方案1 & 方案2 & 方案3 \\
    \midrule
    【指标A】 & 0.82 & 0.76 & 0.91 \\
    【指标B】 & 12.4 & 10.8 & 13.1 \\
    \bottomrule
  \end{tabular}
\end{table}

所有论文图片统一放在 \texttt{figs/} 文件夹。插图示例如\cref{fig:example-result}。
\begin{figure}[htbp]
  \centering
  \includegraphics[width=0.74\textwidth]{example_result.pdf}
  \caption{模型结果示例图（正式提交时替换）}
  \label{fig:example-result}
\end{figure}

【结果分析示例】可以看出，随着【变量】增加，【指标】总体呈【上升/下降】趋势，并在【区间】附近趋于稳定。该规律与【实际机制】一致，说明模型能够较好地描述变量之间的关系。

\subsection{问题二模型的建立与求解}
【示例】在问题一基础上加入资源总量约束。若目标是最小化总成本，可写为
\begin{equation}
  \min_{\bm{x}} F(\bm{x})=\sum_{i=1}^{n} c_i x_i,
  \quad
  \text{s.t.}\quad
  \begin{cases}
    \sum_{i=1}^{n} a_i x_i \geq b,\\
    0\leq x_i\leq u_i,\quad i=1,2,\ldots,n.
  \end{cases}
  \label{eq:q2-model}
\end{equation}
采用【线性规划/整数规划/遗传算法/粒子群算法】求解，得到最优目标值为【结果】，对应决策方案为【方案】。

\subsection{问题三模型的建立与求解}
【示例】针对更复杂情景，在前述模型上引入【随机扰动/多目标权衡/时序变化】。定义综合目标
\begin{equation}
  J=\lambda F_1+(1-\lambda)F_2,
  \label{eq:q3-model}
\end{equation}
其中 $\lambda\in[0,1]$ 表示不同目标之间的权衡系数。通过改变 $\lambda$，得到不同偏好下的候选方案，并结合实际需求给出最终推荐。

% ============================================================
% 七、模型检验与敏感性分析
% ============================================================
\section{模型检验与敏感性分析}

\subsection{模型检验}
若模型属于预测问题，可使用 MAE、RMSE、$R^2$ 或交叉验证评价精度；若属于优化问题，可比较不同算法、不同初值下的结果；若属于机理模型，则应检查结果是否满足实际边界条件。

以均方根误差为例：
\begin{equation}
  \mathrm{RMSE}=\sqrt{\frac{1}{n}\sum_{i=1}^{n}(y_i-\hat y_i)^2}.
\end{equation}
【示例】训练集 RMSE 为【数值】，验证集 RMSE 为【数值】，两者差异较小，模型未出现明显过拟合。

\subsection{敏感性分析}
【示例】选取对模型影响较大的参数 $\theta$，在基准值上下 $10\%$--$20\%$ 范围内扰动并重新计算目标函数。若关键结论与最优方案基本保持不变，则说明模型对该参数不敏感、稳定性较好；反之，应进一步解释敏感原因并给出使用限制。

% ============================================================
% 八、模型评价与推广
% ============================================================
\section{模型评价与推广}

\subsection{模型优点}
\begin{enumerate}
  \item 模型结构与题目各子问题一一对应，变量和约束具有明确的实际含义；
  \item 通过误差检验和敏感性分析验证结果稳定性，而非仅给出单次计算结果；
  \item 求解流程可复现，参数设置、评价指标和输出结果较为清晰。
\end{enumerate}

\subsection{模型不足}
\begin{enumerate}
  \item 部分外部影响因素未被题目数据覆盖，因此模型仍依赖一定的简化假设；
  \item 当样本量较小或数据分布发生明显变化时，模型泛化能力可能下降。
\end{enumerate}

\subsection{模型推广}
【示例】若将【某变量】替换为其他同类指标并重新估计参数，该框架还可用于【相似应用场景】。若获得更丰富的数据，可进一步加入【空间因素/时间因素/不确定性因素】以提升实际适用性。

% ============================================================
% 九、结论
% ============================================================
\section{结论}
本文围绕题目提出的多个子问题完成了数据处理、模型建立、求解与检验。主要结论可概括为：
\begin{enumerate}
  \item 【问题一结论】给出最关键的数值、排序、规律或预测结果；
  \item 【问题二结论】给出最优方案及其主要指标；
  \item 【问题三结论】说明不同情景或参数变化下的最终建议\cite{behzadian2012topsis}；
  \item 【稳健性结论】说明误差、敏感性或验证结果支持模型在什么范围内可靠。
\end{enumerate}

% ============================================================
% AI 工具使用声明：2026 官方要求放在参考文献之前，二选一
% ============================================================
%\CumcmAINotUsed
% 如果竞赛过程中使用了 AI：注释上一行，启用下一行，并据实填写用途；
% 同时编译 support/AI工具使用详情.tex，生成 support/AI工具使用详情.pdf。
\CumcmAIUsed{【请据实填写，如语言润色、代码调试等】}
% 注意：不要把 AI工具使用详情.tex 用 \input 放进 main.tex —— 它是独立文档（自带
% \documentclass 与 \begin{document}），\input 会导致 main.tex 编译失败。
% 正确做法：到 support/ 目录单独编译它生成 PDF：
%     cd support && latexmk AI工具使用详情.tex
% 产出 support/AI工具使用详情.pdf，作为支撑材料提交；正文只需保留上面那句 \CumcmAIUsed 声明。

% ============================================================
% 参考文献：采用 BibTeX 管理（移植自 latex 模板）。
%   1) 真实文献统一维护在根目录 ref.bib（删除下方示例条目前请先添加自己的文献）；
%   2) 正文用 \cite{key} / \citep{key} / \citet{key} 引用，只有被引用到的条目才会列出；
%   3) 如需强制列出 ref.bib 中所有条目，可取消注释：\nocite{*}
%   4) 编译会自动按 xelatex→bibtex→xelatex×2 顺序执行（latexmk / 手动）。
% ============================================================
\bibliographystyle{gbt7714-numerical}
\bibliography{ref}

% ============================================================
% 附录：页数不限；必须包含“支撑材料文件列表”和全部完整源程序
% ============================================================
\CumcmAppendix

\section{支撑材料文件列表}
下面列的是最终支撑材料压缩包中应出现的\textbf{实际文件}，不要只写目录名。正式参赛时请把示例文件替换为真实文件，并保证论文附录代码与支撑材料源码完全一致。
\begin{itemize}
  \SupportFile{code/preprocess.py}{数据清洗、标准化与数据集划分程序；}
  \SupportFile{code/model.py}{模型建立、评价指标与敏感性分析程序；}
  \SupportFile{code/main.py}{完整求解流程入口程序；}
  \SupportFile{requirements.txt}{Python 第三方依赖；}
  \SupportFile{environment.txt}{程序运行环境与运行方式说明；}
  \SupportFile{result/metrics.csv}{模型评价指标结果；}
  \SupportFile{result/sensitivity.csv}{敏感性分析结果；}
  \SupportFile{result/allocation.csv}{优化示例结果。}
  % 若实际使用 AI，启用下面一行（此时须已生成 support/AI工具使用详情.pdf）；未使用 AI 时删除。
  \SupportFile{AI工具使用详情.pdf}{AI 工具使用详情说明。}
\end{itemize}

\section{完整源程序代码}
本节\textbf{直接排入全部实际使用的源程序代码}。不要只摘录“核心代码”。若正式项目包含多个 Python、MATLAB、R、C/C++ 或其他源文件，应逐个使用 \verb|\lstinputlisting| 引入；代码较多时会自然连续排成十几页或几十页，附录页数不受正文 30 页限制。

\subsection{数据预处理程序 \texttt{preprocess.py}}
\lstinputlisting[language=Python,caption={数据预处理完整源程序}]{support/code/preprocess.py}

\subsection{模型与评价程序 \texttt{model.py}}
\lstinputlisting[language=Python,caption={模型与评价完整源程序}]{support/code/model.py}

\subsection{主求解程序 \texttt{main.py}}
\lstinputlisting[language=Python,caption={主求解流程完整源程序}]{support/code/main.py}

% 正式比赛如果还有更多文件，继续逐个加入，例如：
% \subsection{优化程序 \texttt{optimizer.py}}
% \lstinputlisting[language=Python,caption={优化程序完整源程序}]{support/code/optimizer.py}
% \subsection{MATLAB程序 \texttt{solve.m}}
% \lstinputlisting[language=Matlab,caption={MATLAB完整源程序}]{support/code/solve.m}

\end{document}
